% Outline (from the spec):
% - Background: UPI launched 2016, operated by NPCI; Jan Dhan account drive from 2014.
% - Problem: unclear whether adoption raises inclusion or inclusive states adopted faster (direction unclear).
% - Research questions RQ1-RQ4 and hypotheses H1-H4 (below, from the proposal).
\chapter{Introduction}

\section{Background}
% Student writes: UPI launch (2016) and NPCI role; PMJDY launch (2014); policy assumption that
% digital payments widen access to formal finance and support growth; thin state-level evidence.
% Support: \cite{phonepe_pulse} (adoption data), \cite{npci_upi_state} (NPCI state data from Apr 2025),
% \cite{pmjdy_statewise} (PMJDY snapshot), \cite{pradhan2021tfsc} (FI, ICT and growth in states).

\section{Problem statement}
% Student writes: adoption differs across states (internet access, branch density, income, literacy);
% direction unclear (adoption -> inclusion, or inclusion -> adoption); why separating them matters for policy.
% Support: \cite{reddy2024icrier} (diffusion of digital payments), \cite{rbi_fiindex2025} (FI-Index is national only),
% \cite{medianama2025npci}, \cite{dataful2026upi} (state UPI data limits).

\section{Objectives}
\begin{enumerate}
\item To measure how digital payment adoption has grown across Indian states over the available period.
\item To test the relationship between digital payment adoption and financial inclusion.
\item To test the relationship between digital payment adoption and state-level economic growth.
\item To identify the direction of causality (adoption to inclusion, inclusion to adoption, or both).
\item To draw policy implications for lagging states.
\end{enumerate}

\section{Research questions and hypotheses}
\begin{itemize}
\item RQ1: Does higher digital payment adoption raise financial inclusion in a state? (H1: yes, positively.)
\item RQ2: Does adoption contribute to state-level economic growth after controlling for other drivers? (H2: yes, positively.)
\item RQ3: Is the causality bidirectional? (H3: yes; inclusion also encourages adoption.)
\item RQ4: Does the effect differ between better-connected and weaker-infrastructure states? (H4: stronger where internet access is higher.)
\end{itemize}

\section{Structure of the report}
% Student writes: one paragraph per chapter. Part I: Ch. 1 Introduction, Ch. 2 Literature Review.
% Later parts add data (Ch. 3), methods (Ch. 4), results (Ch. 5), conclusion (Ch. 6).
